\section{Two Phases and Three Comings}\label{sec:theology}

The postconciliar Roman calendar opens its account of Advent by assigning the season a twofold character (\textit{duplex indoles}): it prepares the Church for the solemnities of the Nativity, in which the first coming of the Son of God to humanity is recalled, and through that very remembrance it lifts minds and hearts toward Christ's second coming at the end of time, making Advent a time of devout and joyful expectation (\textit{piam et laetam exspectationem}; \textit{Universal Norms on the Liturgical Year and the Calendar}, no.~39).

\subsection{The Double Horizon of the Two Phases}

Because Christ's two visible comings form one mystery of salvation, the postconciliar liturgy does not blur them into a vague sentiment. It separates them in time so that each can shape the believer's imagination in order.

From Evening Prayer~I of the First Sunday of Advent through December~16, the Church turns her face first toward the consummation of history. The Sunday Gospels of the First Week are not Bethlehem narratives; they are Christ's own discourses on the end of the age from Matthew~24, Mark~13, and Luke~21. The weekdays read Isaiah's visions of the mountain of the Lord, the healing of the blind and lame, and the coming of the Lord God with strength, while Preface~I of Advent praises the Father whose Son first assumed the lowliness of human flesh and will come again in glory to open the eternal kingdom to those who watch.

Then, on December~17, the horizon draws abruptly close. Over eight privileged days---December~17 through the morning of December~24---the weekday readings leave the cosmic horizon of the end and step into the courtyard of Abraham, David, Zechariah, Mary, and Joseph. Preface~II of Advent turns directly to the Virgin Mother who waited for her Son with love beyond all telling, and the seven Great O~Antiphons gather the titles of the Old Testament into one nightly plea: \textit{Veni}---``Come!''

Fourth-century catechesis already held these two horizons in one gaze. Teaching the catechumens of Jerusalem on the article of the Creed ``He shall come again in glory,'' St.~Cyril of Jerusalem refused to let Christians treat Bethlehem as a closed chapter, declaring in his fifteenth \textit{Catechetical Lecture} (XV.1):

\begin{quote}
We preach not one advent only of Christ, but a second also, far more glorious than the former. For the former gave a view of His patience; but the latter brings with it the crown of a divine kingdom. \dots\ In His former advent, He was wrapped in swaddling clothes in the manger; in His second, He covereth Himself with light as with a garment. In His first coming, He endured the Cross, despising shame; in His second, He comes attended by a host of Angels, receiving glory. We rest not then upon His first advent only, but look also for His second.
\end{quote}

Cyril immediately adds the practical discipline that guards against both idle speculation and spiritual sleep (\textit{Catechetical Lecture} XV.4): ``venture not thou to declare when these things shall be, nor on the other hand supinely slumber. For he saith, Watch, for in such an hour as ye expect not the Son of Man cometh.'' And when a listener asks how to stand before that coming Judge, Cyril refuses all evasion and quotes the corporal works of mercy of Matthew~25 straight (\textit{Catechetical Lecture} XV.26): ``Learn hence the way; there is here no need of allegory, but to fulfil what is said. \dots\ Be not confident because thou merely possessest the lamp, but constantly keep it burning.''

\subsection{St.~Bernard's Middle Coming}

If Advent looked only backward to Bethlehem and forward to the Last Day, a Christian might be tempted to think Christ absent from the gray December morning in between. St.~Bernard of Clairvaux answers that temptation in his celebrated \textit{Sermons for Advent}, preached to his Cistercian brothers in the twelfth century. After warning in his opening sermon (I.1) that the children of Adam know the name of Advent while remaining too absorbed in perishable things to grasp its meaning, Bernard reveals in his Third and Fifth Sermons that the Lord's coming is not twofold only, but threefold. In his Fifth Sermon for Advent (V.1), Bernard explains the hidden middle advent that joins the manger to the judgment throne:

\begin{quote}
In the first ``He (Christ) was seen on earth and conversed with men,'' when, as He Himself bears witness, ``they both saw and hated'' Him. In the last ``all flesh shall see the Salvation of God,'' and ``they shall look on Him Whom they pierced.'' But in that which intervenes He is hidden, visible only to the elect who see Him in themselves, and so their souls are saved. Therefore, in the first He came in the flesh and in weakness, in the second He comes in spirit and in virtue, and in the third He shall come in glory and majesty. \dots\ Thus, this intermediate advent is the way, so to speak, by which we must travel from the first to the last. In the first Christ was our Redemption, in the third He shall appear as our Life, whilst in this second He is our Repose and our Consolation.
\end{quote}

How does this hidden middle coming take place in an ordinary week? Bernard points straight to the reception of the Word of God and the Eucharist (\textit{Sermons for Advent} V.2--3): ``Keep the word of God in the same way in which the food of thy body can be best preserved; for it is a living bread and the food of thy soul. \dots\ Receive this spiritual bread, therefore, into the stomach of your minds, and having been there properly digested, let it be transmitted thence to the affections and the will.'' Where the Word is heard, kept, and obeyed, the Son comes with the Father to make His abode in the believer (John~14:23), rebuilding ``the Jerusalem of your souls'' and conforming the earthly creature to the image of the heavenly Man.

That middle coming is the very engine of the postconciliar liturgical year. As \textit{Sacrosanctum Concilium} teaches (no.~102), by unfolding the whole mystery of Christ from the Incarnation and Birth to the expectation of His coming, the Church opens to the faithful the riches of her Lord's powers and merits so that these mysteries are in some way made present for all time, and the faithful lay hold of them and are filled with saving grace.

\subsection{Desire as Continual Prayer}

Because Christ comes now in grace to those who wait for Him, Advent is above all the education of longing. St.~Augustine gave the classic account of that longing in his \textit{Exposition on Psalm~37} (Vulgate numbering; Psalm~38 in Hebrew, \S13), reflecting on St.~Paul's command to ``pray without ceasing'' (1~Thessalonians~5:17, proclaimed on the Third Sunday of Advent in Year~B):

\begin{quote}
For it is thy heart's desire that is thy prayer; and if thy desire continues uninterrupted, thy prayer continueth also. \dots\ There is another inward kind of prayer without ceasing, which is the desire of the heart. Whatever else you are doing, if you do but long for that Sabbath, you do not cease to pray. If you would never cease to pray, never cease to long after it. The continuance of thy longing is the continuance of thy prayer.
\end{quote}

Writing to the Roman noblewoman Proba (\textit{Letter} 130.18), Augustine explains why fixed hours of vocal prayer remain indispensable even when prayer is rooted in desire: ``at certain stated hours and seasons we also use words in prayer to God, that by these signs of things we may admonish ourselves, and may acquaint ourselves with the measure of progress which we have made in this desire, and may more warmly excite ourselves to obtain an increase of its strength.'' Spoken psalms, hymns, and antiphons do not inform God of what He does not know; they stretch our narrow hearts until they are wide enough to receive what He longs to give.

\subsection{Mary, Model of the Expectant Church}

No human being ever waited for Christ as the Virgin Mary did. The Second Vatican Council placed Mary squarely at the heart of the Church's liturgy and mystery (\textit{Lumen Gentium}, nos.~52--56): redeemed in a more exalted manner by reason of the merits of her Son, she is greeted by the Fathers as the New Eve who, believing the angel's word without doubt, untied the knot of Eve's disobedience and gave Life itself to the world. In the postconciliar Advent, Mary's presence grows steadily brighter: celebrated as the Immaculate Virgin on December~8, honored across the Americas as Our Lady of Guadalupe on December~12, singing her \textit{Magnificat} as the Responsorial Canticle on the Third Sunday of Year~B, and standing at the center of the Annunciation and Visitation Gospels on the Fourth Sunday of Years~B and~C and throughout December~17--24.

Walking beside her keeps Advent from becoming grim. The season wears violet, omits the \textit{Gloria}, and moderates flowers and music, yet never drops the \textit{Alleluia}. That balance captures the exact spirit of \textit{Universal Norms} no.~39: devout and joyful expectation (\textit{piam et laetam exspectationem}). A people that knows its own sins cannot be flippant before the Judge of the living and the dead; a people that knows the Redeemer has already taken flesh from Mary cannot despair. Between sober repentance and joyful hope, Advent teaches the heart to stand awake until the Daystar rises.
